\scanpage{115}
\mspage{121}
b) Pour tout $F\in\operatorname{Ob}\unsure{\mathbb D|D}(\unsure{x|X})$ \uline{loc.} \uline{constant}, $f_*F$ est loc. constant sur $Y$ et se calcule ``fibre par fibre'' -- et \unsure{itou|idem} après tout changement de base

c) Pour tout $Y''\underset{p}{\unsure{\mathord{-}|\to}}Y'$ sur $Y$, si $p$ est une $\mathbb D$-équiv. cohomologique, alors il en est de même de $q:X''\to X'$.

c$'$) Comme c, mais on suppose \deleted{que} comme hyp. sur $p$ que $Y''$ soit la cat. finale, et $Y'$ ait un objet final.

\uline{Dém} Dans l'ordre ($\pm$ \unsure{tautologique})

\[a\Longrightarrow b\Longrightarrow c\deleted{\unsure{s|'}}\Longrightarrow c'\Longrightarrow a\]

\uline{Définition 5}. Dans les conditions équivalentes de la prop. 9, on dit que $f$ est une $\mathbb D$-\uline{fibration} \uline{cohomologique}. Sous les conditions duales, on dit que $f$ est une $\mathbb D$-\uline{fibration homologique}. On dit que $f$ est une $\mathbb D$-\uline{fibration forte}, si c'est à la fois une $\mathbb D$-fibration cohomologique et une $\mathbb D$-fibration homologique.

Ces notions \deleted{\unsure{toutes deux}} renforcent la notion de $\mathbb D$-fibration. \inserted{Elles ont \unsure{l'avantage} de plus d'être \uline{stables par composition}, \unsure{ainsi que|car il suffit} \unsure{par|pour} chgt. de base.}

\deleted{\unsure{La} \illegible{2} à l'aide de ces deux \unsure{notions} le résultat \unsure{suivant} :}

\deleted{\uline{Prop. 10} Soit $f:X\to Y$ dans \unsure{Diag}. Conditions équivalentes :}

\deleted{a) $\forall\, y\in\operatorname{Ob}Y$, $X_y-X_{/y}$ est une $\mathbb D$-équiv. $1$-complète (resp. $2$ complète), et de $=$ pour $X_{/y}-X_{/y'}$ pour toute flèche $y-y'$ dans $Y$. Et ça reste valable après tout chgt de base (NB il suffit de l'exiger pour la première de ces conditions, concernant $X_y-X_{/y}$.)}

\deleted{\unsure{a$'$|b}) Condition duale de a, en termes des $X_y-{}_{y\backslash}X$, $\,{}_{y'\backslash}X\to{}_{y\backslash}X$, qui doivent être \unsure{tous|t.} des $\mathbb D$-équiv. $1$-complètes (resp. $2$-complètes), et ça reste vrai après tt chgt. de base.}

\deleted{\unsure{b)}}
% archive sheet number: 115
